\subsection{Confidence Sets, Pivots, and Asymptotic Intervals}

Confidence procedures attach uncertainty to estimators by controlling coverage probability under repeated sampling.

\begin{definition} \label{def:confidence-set-coverage} \hlt{Confidence Set}\\
A random set $C(X)\subseteq\Psi$ is a \hlt{$1-\alpha$ confidence set} for $\psi(\theta)$ if
\[
\inf_{\theta\in\Theta}\mathbb{P}_\theta\bigl(\psi(\theta)\in C(X)\bigr)\geq1-\alpha.
\]
The infimum is the \hlt{coverage guarantee}.
\end{definition}

\begin{definition} \label{def:pivot-confidence-construction} \hlt{Pivot}\\
A \hlt{pivot} is a function $Q(X,\theta)$ whose distribution under $P_\theta$ does not depend on $\theta$.  If $a_\alpha\leq Q(X,\theta)\leq b_\alpha$ has probability $1-\alpha$, solving these inequalities for $\theta$ gives an exact confidence set.
\end{definition}

\begin{example} \label{ex:normal-mean-known-variance-confidence-interval} \hlt{Normal Mean with Known Variance}\\
If $X_1,\ldots,X_n$ are independent $\mathcal{N}(\mu,\sigma^2)$ observations with known $\sigma$, then
\[
Z=\frac{\sqrt{n}(\overline{X}-\mu)}{\sigma}\sim\mathcal{N}(0,1).
\]
Thus
\[
\left[\overline{X}-z_{1-\alpha/2}\frac{\sigma}{\sqrt{n}},
\overline{X}+z_{1-\alpha/2}\frac{\sigma}{\sqrt{n}}\right]
\]
has exact coverage $1-\alpha$.
\end{example}

\begin{proposition} \label{prop:test-confidence-duality} \hlt{Test--Confidence Duality}\\
For each $\psi_0\in\Psi$, let $\varphi_{\psi_0}$ be a level-$\alpha$ test of $H_0:\psi(\theta)=\psi_0$.  Then
\[
C(X)=\{\psi_0:\varphi_{\psi_0}(X)=0\}
\]
is a $1-\alpha$ confidence set.
\begin{proof}
If the true value is $\psi(\theta)$, then the event $\psi(\theta)\notin C(X)$ is the event that the level-$\alpha$ test of the true null rejects.  Its probability is at most $\alpha$.
\end{proof}
\end{proposition}

\begin{definition} \label{def:wald-score-likelihood-intervals} \hlt{Wald, Score, and Likelihood Intervals}\\
For a scalar regular parameter, three common asymptotic confidence constructions are:
\begin{enumerate}[label=(\roman*)]
\setlength{\itemsep}{0pt}
\item \hlt{Wald}: center at $\widehat{\theta}$ and use estimated standard error from observed or expected information.
\item \hlt{Score}: invert score tests under the null value.
\item \hlt{Likelihood}: invert likelihood-ratio tests using the approximation $-2\log\lambda(X)\Rightarrow\chi^2_1$.
\end{enumerate}
\hlt{Key Insight}: These agree to first order in regular large samples but may differ substantially in small or boundary-sensitive problems.
\end{definition}
